\section*{Referências}
\addcontentsline{toc}{section}{Referências}
\small\singlespacing
\begin{enumerate}[label={[\arabic*]},leftmargin=1.1cm,itemsep=0.35em]
\item PONTIFÍCIA UNIVERSIDADE CATÓLICA DE MINAS GERAIS. \textit{XIII -- Estratégias de crescimento, viabilidade e coerência}. Orientação da atividade, Tópicos III, 2026.
\item KASSHINOKUN. \textit{Q3\_Q4\_2026\_Public: documentação do aplicativo PetCuida}. GitHub. Disponível em: \url{https://github.com/kasshinokun/Q3_Q4_2026_Public/tree/main/TOPICOS_III}. Acesso em: 26 out. 2026.
\item CASSINO, G. S.; RIBEIRO, G. N.; MELO, J. R. V.; BARBOSA, K. R. L. \textit{Relatório de Testes de Usabilidade -- PetCuida Alpha 3i}. Atividade VI, 2026. Disponível em: \url{https://github.com/kasshinokun/Q3_Q4_2026_Public/tree/main/TOPICOS_III/Atividade_VI}. Acesso em: 26 out. 2026.
\item CAMUFFO, A.; CORDOVA, A.; GAMBARDELLA, A.; SPINA, C. A scientific approach to entrepreneurial decision making: evidence from a randomized control trial. \textit{Management Science}, v. 66, n. 2, p. 564--586, 2020. DOI: 10.1287/mnsc.2018.3249.
\item RIES, E. \textit{The lean startup}. New York: Crown Business, 2011.
\item BLANK, S. Why the lean start-up changes everything. \textit{Harvard Business Review}, v. 91, n. 5, p. 63--72, 2013.
\item EISENMANN, T.; RIES, E.; DILLARD, S. \textit{Hypothesis-driven entrepreneurship: the lean startup}. Boston: Harvard Business School, 2012. (Background Note 812-095).
\item MARMER, M.; HERRMANN, B. L.; DOGRULTAN, E.; BERMAN, R. \textit{Startup Genome Report Extra on Premature Scaling}. Startup Genome, 2011.
\item ROCHET, J.-C.; TIROLE, J. Platform competition in two-sided markets. \textit{Journal of the European Economic Association}, v. 1, n. 4, p. 990--1029, 2003. DOI: 10.1162/154247603322493212.
\item PARKER, G. G.; VAN ALSTYNE, M. W.; CHOUDARY, S. P. \textit{Platform revolution}. New York: W. W. Norton, 2016.
\item ANSOFF, H. I. Strategies for diversification. \textit{Harvard Business Review}, v. 35, n. 5, p. 113--124, 1957.
\item GUPTA, S.; LEHMANN, D. R.; STUART, J. A. Valuing customers. \textit{Journal of Marketing Research}, v. 41, n. 1, p. 7--18, 2004. DOI: 10.1509/jmkr.41.1.7.25084.
\item REICHHELD, F. F. The one number you need to grow. \textit{Harvard Business Review}, v. 81, n. 12, p. 46--54, 2003.
\item WEINBERG, G.; MARES, J. \textit{Traction: how any startup can achieve explosive customer growth}. New York: Portfolio, 2015.
\item ELLIS, S.; BROWN, M. \textit{Hacking growth}. New York: Crown Business, 2017.
\item CROLL, A.; YOSKOVITZ, B. \textit{Lean analytics}. Sebastopol: O'Reilly, 2013.
\item NIELSEN, J.; LANDAUER, T. K. A mathematical model of the finding of usability problems. In: \textit{Proceedings of INTERCHI '93}. New York: ACM, 1993. p. 206--213. DOI: 10.1145/169059.169166.
\item BROOKE, J. SUS: a `quick and dirty' usability scale. In: JORDAN, P. W. et al. (ed.). \textit{Usability evaluation in industry}. London: Taylor \& Francis, 1996. p. 189--194.
\item ABEMPET. \textit{Balanço Abempet 2025: setor pet faturou R\$ 77,96 bilhões em 2025}. Câmara Setorial de Animais de Estimação, MAPA, 2026. Disponível em: \url{https://www.gov.br/agricultura/pt-br/assuntos/camaras-setoriais-tematicas/documentos/camaras-setoriais/animais-e-estimacao/2026/45a-ro-23-04-2026/release-projecao-2025-setor-pet.pdf}. Acesso em: 26 out. 2026.
\item ABINPET; ABEMPET. \textit{Dados de mercado do setor pet}. 44ª Reunião Ordinária da Câmara Setorial, MAPA, 30 out. 2025. Disponível em: \url{https://www.gov.br/agricultura/pt-br/assuntos/camaras-setoriais-tematicas/documentos/camaras-setoriais/animais-e-estimacao/2025/44a-ro-30-10-2025/release-projecao-2025-setor-pet.pdf}. Acesso em: 26 out. 2026.
\item STERN, J. Censo pet: Brasil ultrapassa 164 milhões, com destaque para gatos e pequenos animais. \textit{Petfood Forum Brasil}, 15 jul. 2025. Disponível em: \url{https://www.petfoodforumbrasil.com/ultimas-noticias/article/15750669/censo-pet-brasil-ultrapassa-164-milhes-com-destaque-para-gatos-e-pequenos-animais}. Acesso em: 26 out. 2026.
\item CNN BRASIL. Brasil é 3º país com mais pets e setor prevê faturar R\$ 81 bilhões em 2026. 23 set. 2026. Disponível em: \url{https://www.cnnbrasil.com.br/economia/money/negocios/brasil-e-3o-pais-com-mais-pets-e-setor-preve-faturar-r-81-bilhoes-em-2026/}. Acesso em: 26 out. 2026.
\item LAVALLEE, E.; MUELLER, M. K.; McCOBB, E. A systematic review of the literature addressing veterinary care for underserved communities. \textit{Journal of Applied Animal Welfare Science}, v. 20, n. 4, p. 381--394, 2017. DOI: 10.1080/10888705.2017.1337515.
\item ACCESS TO VETERINARY CARE COALITION. \textit{Access to veterinary care: barriers, current practices, and public policy}. Knoxville: University of Tennessee, 2018.
\item CONSELHO FEDERAL DE MEDICINA VETERINÁRIA. \textit{Resolução nº 1.465, de 13 de julho de 2022}. Dispõe sobre a telemedicina veterinária. Brasília: CFMV, 2022.
\item BRASIL. \textit{Lei nº 13.709, de 14 de agosto de 2018}. Lei Geral de Proteção de Dados Pessoais (LGPD). Brasília: Presidência da República, 2018.
\item HOFFMAN, R.; YEH, C. \textit{Blitzscaling}. New York: Currency, 2018.
\item AUTORIDADE NACIONAL DE PROTEÇÃO DE DADOS. \textit{Guia orientativo sobre segurança da informação para agentes de tratamento de pequeno porte}. Brasília: ANPD, 2021.
\item CASSINO, G. S.; RIBEIRO, G. N.; MELO, J. R. V.; BARBOSA, K. R. L. \textit{Atividade X -- Tipo societário e aspectos legais: PetCuida}. 2026. Disponível em: \url{https://github.com/kasshinokun/Q3_Q4_2026_Public/tree/main/TOPICOS_III/Atividade_X}. Acesso em: 26 out. 2026.
\item CASSINO, G. S.; RIBEIRO, G. N.; MELO, J. R. V.; BARBOSA, K. R. L. \textit{Atividade XI -- Estrutura organizacional com descrição de responsabilidades}. Versão final, 2026.
\end{enumerate}
\normalsize\onehalfspacing
